\section{PIC 初始化：ICW 序列详解}

PIC 通过 4 个 ICW（Initialization Command Word）完成初始化。
每个 ICW 有精确的格式和发送顺序：

\codefile{src/kernel/arch/pic.c}
\begin{lstlisting}[style=cstyle]
void pic_init(void) {
    // 保存当前 mask（可选，我们会覆盖）
    uint8_t a1 = inb(PIC1_DATA);
    uint8_t a2 = inb(PIC2_DATA);
\end{lstlisting}

\subsection{ICW1：开始初始化}

\codefile{src/kernel/arch/pic.c}
\begin{lstlisting}[style=cstyle]
    outb(PIC1_CMD, 0x11);  // ICW1: INIT + ICW4_NEEDED
    io_wait();
    outb(PIC2_CMD, 0x11);
    io_wait();
\end{lstlisting}

\hex{11} = ICW1\_INIT (bit 4=1) | ICW1\_ICW4 (bit 0=1)。

向命令端口写入 ICW1 告诉 PIC："开始初始化序列，后面会有 ICW4"。

\begin{whybox}
\texttt{io\_wait()} 是什么？ISA 总线的某些旧设备需要在连续 I/O 操作间
插入一个短延迟。常见实现是向端口 \hex{80}（POST 诊断端口）写一个字节：

\codefile{src/include/arch/io.h}
\begin{lstlisting}[style=cstyle]
static inline void io_wait(void) {
    outb(0x80, 0);  // 约 1 微秒延迟
}
\end{lstlisting}
\end{whybox}

\subsection{ICW2：向量偏移}

\codefile{src/kernel/arch/pic.c}
\begin{lstlisting}[style=cstyle]
    outb(PIC1_DATA, 0x20);  // master: IRQ0 → 向量 0x20
    io_wait();
    outb(PIC2_DATA, 0x28);  // slave:  IRQ8 → 向量 0x28
    io_wait();
\end{lstlisting}

这就是\textbf{remap 的核心}：告诉每片 PIC 它管理的 IRQ 对应的起始向量号。

Remap 后的映射：

\begin{longtable}{ccl}
  \toprule
  \textbf{IRQ} & \textbf{向量} & \textbf{设备} \\
  \midrule
  0  & \hex{20} (32) & PIT Timer \\
  1  & \hex{21} (33) & PS/2 Keyboard \\
  2  & \hex{22} (34) & Cascade（slave PIC） \\
  8  & \hex{28} (40) & CMOS RTC \\
  14 & \hex{2E} (46) & Primary ATA \\
  \bottomrule
\end{longtable}

\subsection{ICW3：级联关系}

\codefile{src/kernel/arch/pic.c}
\begin{lstlisting}[style=cstyle]
    outb(PIC1_DATA, 0x04);  // master: slave 挂在 IRQ2 (bit 2 = 1)
    io_wait();
    outb(PIC2_DATA, 0x02);  // slave: cascade identity = 2
    io_wait();
\end{lstlisting}

Master 的 ICW3 是一个 bitmask：bit $N$ = 1 表示 IRQ$N$ 上挂了一个 slave PIC。
\hex{04} = bit 2 = IRQ2 上有 slave。

Slave 的 ICW3 是自身的 cascade identity：2 表示"我挂在 master 的 IRQ2 上"。

\begin{figure}[H]
  \centering
  % figures/ch05/fig02-master-pic-icw3-04-irq2-1.tex — auto-extracted from chapter source
\begin{tikzpicture}[node distance=0cm]
  \foreach \i/\label in {0/0, 1/0, 2/\textbf{1}, 3/0, 4/0, 5/0, 6/0, 7/0} {
    \node[bitcell, fill=white] (m\i) at (\i*0.8, 0) {\label};
    \node[font=\tiny\ttfamily, above=0.1cm of m\i] {IRQ\i};
  }
  \node[below=0.5cm of m2, font=\small\itshape] {Slave 挂在 IRQ2};
  \draw[pointer] (m2.south) -- ++(0, -0.3);
\end{tikzpicture}

  \caption{Master PIC 的 ICW3 = \hex{04}：IRQ2 位置为 1}
\end{figure}

\subsection{ICW4：模式设置}

\codefile{src/kernel/arch/pic.c}
\begin{lstlisting}[style=cstyle]
    outb(PIC1_DATA, 0x01);  // 8086/88 模式
    io_wait();
    outb(PIC2_DATA, 0x01);
    io_wait();
\end{lstlisting}

\hex{01} = ICW4\_8086：告诉 PIC 工作在 8086 模式（相对于 8085 模式）。
现代 x86 都用这个。

\subsection{初始 Mask：全部关闭}

\codefile{src/kernel/arch/pic.c}
\begin{lstlisting}[style=cstyle]
    outb(PIC1_DATA, 0xFF);  // mask 所有 master IRQ
    outb(PIC2_DATA, 0xFF);  // mask 所有 slave IRQ
}
\end{lstlisting}

\begin{designbox}
初始化后先把所有 IRQ mask 掉。原因：
\begin{enumerate}
  \item 此时大部分 IRQ 还没安装 handler。如果中断进来但 IDT 中没有有效的 gate，CPU 会触发 \#GP。
  \item 每个驱动在 init 时自己调用 \texttt{pic\_clear\_mask(irq)} 打开需要的 IRQ — "谁用谁开"原则。
\end{enumerate}
\end{designbox}

% ============================================================================

